\section{Related Work}
In a MAD system, each LLM agent presents its viewpoint and scrutinizes the viewpoints of other LLM agents across multiple rounds of debate~\citep{sun2024corex}. In summary, the existing MAD methods can be categorized as two types, namely \textit{sequential debate} and \textit{parallel debate}.

\textbf{Sequential Debate.} In these methods~\citep{hu2025debate, brownscalable, michael2023debate, wang2025learning, he2024agentscourt}, LLM agents generate their viewpoints in turn. Each LLM agent can only obtain the viewpoints of its preceding LLM agents. For example, \citet{liang2023encouraging} require two LLMs to refute each other in turn. In addition to debaters, \citet{guan2025mmd} add extra roles, such as judge and critic. The judge speaks before debaters to explain the task, and the critic speaks last to summarize debates. However, in a sequential debate system, each LLM agent must wait for previous LLM agents to finish reasoning before it starts. This makes debating time increase linearly with the number of LLM agents, leading to low efficiency which limits the scalability.

\textbf{Parallel Debate.} In these methods~\citep{phamlet, yin2023exchange, chern2024can, khandebating, liang2024debatrix, li2024coevol, hegazy2024diversity, zhang2024breaking}, all LLM agents simultaneously generate their viewpoints based on the viewpoints of other LLM agents in the last debating round. For example, \citet{chanchateval} require LLM agents to critique all answers generated in the last debating round and update its answer in each debating round simultaneously. In addition to the answers generated in the last round, \citet{sun2024towards} also provide each LLM agent with task-related information retrieved from the web. Besides, some methods~\citep{duan2024enhancing, yoffe2024debunc} try to adjust the debating influence of each LLM agent to improve the debating effectiveness. For example, \citet{chen2023reconcile} require each LLM agent to generate the confidence score for its own answer, and then inputs the score to other LLM agents along with the answer.

Since sequential debate systems face the low-efficiency issue mentioned above, our proposed CortexDebate follows the parallel debate framework. Compared with existing parallel debating methods which require each LLM agent to debate with all others in each round, our CortexDebate dynamically decides the necessary debating agents by establishing a sparse debating graph among all involved LLM agents, so that the input context to each agent can be shortened. This is also in contrary to~\citep{liu2024groupdebate, li2024improving} in which the debating opponents are fixed. Furthermore, different from prior methods which determine the debating impact of each LLM agent simply based on its own confidence, we introduce the McKinsey Trust Formula so that both the confidence of each LLM agent and the usefulness to its debating component can be evaluated.
